\section{Kết luận}
\label{sec:conclusion}

Báo cáo này đã trình bày tổng quan về bài toán jailbreaking trong các mô hình
ngôn ngữ lớn, bao gồm nền tảng kỹ thuật, các dạng tấn công phổ biến, các hướng
phòng thủ, phương pháp đánh giá, benchmark liên quan và rủi ro mở rộng trong
các hệ thống agentic. Thông qua các phần tổng hợp lý thuyết, báo cáo cho thấy
jailbreaking không chỉ là vấn đề của một prompt đơn lẻ, mà là một bài toán hệ
thống liên quan đến instruction hierarchy, alignment, threat model, guardrail,
evaluation và trade-off giữa safety và utility.

Về mặt thực nghiệm, nhóm chọn hướng tự xây dựng demo từ đầu theo yêu cầu của
đề tài. Demo không nhằm tái lập đầy đủ một phương pháp cụ thể như GCG, PAIR hay
AutoDAN, mà tập trung cài đặt một pipeline nhỏ có kiểm soát gồm các thành phần
chính được đề cập trong tutorial: attack transformations, defense
configurations, evaluation metrics và visualization. Pipeline được áp dụng trên
một tập prompt nhỏ gồm các nhóm \texttt{knowledge}, \texttt{company} và
\texttt{hazardous}, sau đó đánh giá qua nhiều attack type và defense type khác
nhau.

Kết quả thực nghiệm cho thấy khi không có defense bổ sung, hệ thống vẫn tồn tại
unsafe compliance với ASR bằng 0.133. Trong các defense được thử nghiệm,
\texttt{output\_moderation} đạt trade-off tốt nhất: ASR giảm về 0.000 trong
khi Utility Rate vẫn giữ ở mức 0.613. Ngược lại, \texttt{input\_filter} và
\texttt{defense\_in\_depth} cũng giúp giảm ASR nhưng làm tăng over-refusal,
đặc biệt ở nhóm \texttt{hazardous}. Điều này cho thấy một defense tốt không chỉ
cần giảm unsafe compliance, mà còn phải hạn chế việc từ chối nhầm các yêu cầu
hợp lệ.

Một kết quả đáng chú ý khác là rủi ro không phân bố đều giữa các nhóm prompt.
Nhóm \texttt{company} có ASR cao nhất khi không có defense, cho thấy các yêu
cầu liên quan đến thông tin nội bộ, secret hoặc instruction hierarchy là nhóm
cần được kiểm soát cẩn thận. Trong khi đó, nhóm \texttt{knowledge} không tạo ASR
trong thí nghiệm này nhưng vẫn có over-refusal, phản ánh khó khăn của hệ thống
trong việc phân biệt giữa thảo luận học thuật về jailbreak và yêu cầu gây hại
thực sự.

Bên cạnh experiment chính, nhóm cũng thực hiện manual web evaluation như một
phân tích bổ sung. Phần này tổng hợp các nhãn \texttt{Passed\_*} từ dữ liệu
manual và không được gộp vào các metric chính như ASR, ORR hoặc Utility Rate.
Kết quả manual giúp xác định các case cần đọc sâu hơn, nhưng không được xem là
bằng chứng định lượng thay thế cho pipeline đánh giá chính.

Do hạn chế về tài nguyên tính toán, quota API và thời gian thực hiện, thí
nghiệm trong báo cáo có phạm vi giới hạn. Bộ prompt còn nhỏ, target model chính
được gọi thông qua OpenRouter API, evaluator hiện tại là rule-based, và các
defense mới chỉ ở mức minh họa. Vì vậy, các kết quả không nên được hiểu như kết
luận tổng quát về mọi mô hình hoặc mọi dạng jailbreak. Giá trị chính của demo
nằm ở việc minh họa được quy trình đánh giá, cách đo lường safety--utility
trade-off, và các dạng lỗi thường gặp như unsafe compliance và over-refusal.

Trong tương lai, có thể mở rộng công việc theo nhiều hướng. Thứ nhất, dataset
cần được mở rộng bằng các benchmark công khai và nhiều nhóm rủi ro hơn. Thứ
hai, evaluator nên được cải thiện bằng human evaluation, LLM-as-a-judge hoặc
guard model chuyên dụng. Thứ ba, pipeline nên được chạy trên nhiều target model
để so sánh độ bền của từng mô hình. Thứ tư, có thể tích hợp các defense mạnh
hơn được đề cập trong tutorial, chẳng hạn smoothing, adversarial training,
circuit breaker hoặc contextual moderation. Cuối cùng, hướng mở rộng quan trọng
là đánh giá jailbreak trong agentic systems, nơi prompt injection có thể ảnh
hưởng đến tool-use, memory, retrieval và workflow nhiều bước.

Tổng kết lại, báo cáo cho thấy jailbreaking là một bài toán đánh giá an toàn
phức tạp, cần được tiếp cận theo hướng hệ thống. Một mô hình hoặc ứng dụng LLM
không thể được xem là an toàn chỉ vì có tỷ lệ từ chối cao, cũng không thể được
xem là hữu ích nếu bỏ qua rủi ro unsafe compliance. Đánh giá jailbreak cần cân
bằng giữa ba mục tiêu: giảm khả năng làm theo yêu cầu nguy hiểm, hạn chế
over-refusal và duy trì utility cho các yêu cầu hợp lệ.